%==============================================================================
% UQFF Manuscript v2 — Section 3
% Topic:    Closure architecture (5 namespaces, ~1080 named closures, software shape)
% Status:   first draft, ready for Daniel's review
% Anchors:  CLOSURE_ATLAS.md, ARCHITECTURE.md, calculate_status_report()
%==============================================================================

\section{Closure architecture}
\label{sec:architecture}

A \emph{closure} in this framework is a named numerical identity
derived from the nine primitives of Sec.~\ref{sec:primitives} and
yielding a value comparable to a measured physical quantity. This
section describes how the closures are organized in the
implementation, so that a reader who wants to inspect any
particular closure cited elsewhere in this manuscript can locate
the underlying code quickly.

\subsection{The five closure namespaces}
\label{sec:five_namespaces}

The calculator's closures are distributed across five
distinct dispatch namespaces. The split reflects the developmental
history of the framework rather than any conceptual hierarchy; all
five draw on the same nine primitives.

\begin{table}[H]
\centering
\caption{The five closure namespaces, with current dispatch-key
counts and the corresponding canonical access pattern. Counts are
the live values reported by the calculator at the time of this
writing; the fidelity gate of Sec.~\ref{sec:fidelity_gate} asserts
the structural integrity of each on every commit.}
\label{tab:five_namespaces}
\begin{tabular}{l r p{6.5cm}}
\toprule
\textbf{Namespace} & \textbf{Count} & \textbf{Access pattern + use} \\
\midrule
\texttt{PARADOX\_TO\_CLOSURE} & 794 & \texttt{calculate\_paradox(\{"paradox":NAME\})}; the broadest dispatch, covers most named identities (Yang-Mills, magic numbers, Holmlid, etc.) \\
\texttt{PARADOX\_TO\_MILLENNIUM} & 8 & \texttt{calculate\_paradox} aliases for the seven Clay Millennium prize problems plus the black-hole information paradox \\
\texttt{calculate\_lenr\_full} sub-keys & 10 & Holmlid, Parkhomov, Mizuno, Pons-Fleischmann, Rossi Early/X/SK, Star-Magic reactor, Widom-Larsen, Kozima TNCF, meson cascade \\
\texttt{calculate\_nuclear\_magic} sub-keys & 7 & Magic numbers + BE/$A$ peak + $\alpha$-particle binding + deuteron binding + Coulomb 626\,eV cross-check + Mayer-Jensen comparison \\
Bucket observables (9 surfaces) & 248 & Cosmology, particle physics, GW events, AGN jets, astrophysics, high-energy astro, QGP, Higgs precision, BSM constraints --- each surface returns an \texttt{observables} list \\
\midrule
\textbf{Total named closures} & \textbf{$\sim 1{,}067$} & Plus 530 legacy-freeform entries; gross 793 schema-tagged after deduplication \\
\bottomrule
\end{tabular}
\end{table}

The bucket-observable count (248) is itself the sum of nine bucket
surfaces returning lists of observables (cosmology 18, particle
physics 49, GW events 20, AGN jets 20, astrophysics 20, high-energy
astro 7, QGP 4, Higgs precision 8, BSM constraints 9, plus a few
additional rows in each). The five-namespace decomposition is
canonicalized by the \texttt{uqff\_cli.py} dispatcher, which walks
all five in order during \texttt{uqff predict NAME} and
\texttt{uqff search SUBSTR} queries.

\subsection{Backing whitepaper corpus}
\label{sec:whitepaper_corpus}

Every closure in the namespaces of Table~\ref{tab:five_namespaces}
has a \texttt{primary\_source} attribute pointing at a backing
whitepaper (PAPER\_XXXX), with current corpus size 1,994 papers
covering both schema-tagged and legacy-freeform closures. The
package's \texttt{uqff proofs} subcommand exposes the entire corpus:

\begin{verbatim}
$ uqff proofs list --filter holmlid
PAPER_1133_Holmlid_Rydberg_SCm_Bridge.md
PAPER_1136_Holmlid_KER_Reactor_Validation.md
PAPER_1137_Holmlid_Rossi_Parkhomov_Validation.md
PAPER_1138_Holmlid_Parkhomov_Pons_Fleischmann_Upgrade.md
... (additional matches)
$ uqff proofs show PAPER_1133
[full whitepaper contents]
\end{verbatim}

The whitepaper-to-closure mapping is audited in
\texttt{COVERAGE\_GAPS.md} (bundled), which reports both
the wired-coverage rate (closures with a backing whitepaper) and
the inverse-coverage rate (whitepapers referenced by at least one
closure). These two rates do not have to be 100\,\% --- some
whitepapers are historical scaffolding; some closures span multiple
papers --- but both are monitored for drift on every commit.

\subsection{Master inventory document}
\label{sec:master_inventory}

The master inventory of all closures, observables, and their
canonical sources is \texttt{CLOSURE\_ATLAS.md}, generated and
maintained by the package's own dispatch-introspection machinery.
At the time of this writing, the atlas reports approximately
4{,}164 distinct closure-related artifacts (named closures, backing
whitepapers, chain-trace audit files, arXiv-bundle papers, and
Lean 4 formal-scaffold theorems). Readers who want to confirm any
specific claim in the manuscript can locate the underlying code,
the backing whitepaper, the chain-trace audit log, and the
formal-scaffold theorem statement from a single document.

\subsection{Software architecture brief}
\label{sec:software_architecture}

The framework's runtime architecture is intentionally minimal:

\begin{itemize}
  \item One file (\texttt{uqff\_pure\_calculator.py}, $\sim$2.7\,MB,
        $\sim$48\,000 lines) holds the entire calculator.
  \item Thirty-four public surfaces (\texttt{calculate\_*} functions)
        return uniformly-shaped \texttt{\{'value': X\}} dicts.
  \item Zero runtime dependencies (the calculator imports only the
        Python standard library).
  \item Five auxiliary files provide CLI (\texttt{uqff\_cli.py}),
        REST API (\texttt{uqff\_api.py}), Jupyter integration
        (\texttt{uqff\_jupyter.py}), test gate
        (\texttt{uqff\_fidelity\_tests.py}), and the C++ reference
        (\texttt{uqff\_exact\_closures.cpp}).
\end{itemize}

The single-file calculator design is deliberate. The framework's
audit history (Tier-3 G8/G9/G10/G11 audits, documented in
\texttt{SECURITY\_AND\_MEMORY\_AUDIT.md} and
\texttt{PERFORMANCE\_PROFILE.md}) shows the design trade-off as net
favorable: cold-import time is 537\,ms (modest), per-dispatch latency
is 3.4\,$\mu$s (290{,}000 calls per second on a single thread),
memory footprint is 6.8\,MB module + 224\,MB process RSS, and the
single-file structure simplifies the cross-language verification
(Sec.~\ref{sec:cpp_crosscheck}).

\subsection{Reader's how-to-locate guide}
\label{sec:how_to_locate}

For any closure value cited elsewhere in this manuscript:

\begin{enumerate}
  \item Open a Python shell after \texttt{pip install uqff}.
  \item Run \texttt{uqff search NAME} to find the closure across
        all five namespaces.
  \item Run \texttt{uqff predict NAME} to retrieve the value plus
        the canonical-source whitepaper attribution.
  \item Run \texttt{uqff proofs show PAPER\_XXXX} to read the
        backing whitepaper.
  \item For the calculator source, the function name is the
        closure name prefixed with \texttt{\_paradox\_proof} or
        appears as a key in the dispatch dict; the file is
        \texttt{uqff\_pure\_calculator.py}, line $\sim$200 (the
        \texttt{PARADOX\_TO\_CLOSURE} dict starts here) onwards.
\end{enumerate}

There is no closure in this manuscript whose underlying code,
backing whitepaper, and cross-language C++ counterpart cannot be
located within five minutes by a reviewer using only the
\texttt{uqff} CLI and a text editor.

